\section{Metodologia del modelo de inventario}

La simulacion se implemento por periodos diarios, usando un horizonte de $365$ dias y $10$ replicas independientes por politica. La semilla base fue $20260626$ y las semillas de cada replica se derivaron de esa base para conservar reproducibilidad.

Los parametros comunes fueron:

\begin{itemize}
    \item Demanda diaria: Poisson con media de $20$ unidades por dia.
    \item Inventario inicial: $120$ unidades.
    \item Tiempo de entrega: $3$ dias.
    \item Costo fijo de orden: $120$.
    \item Costo de mantenimiento: $0.35$ por unidad por dia.
    \item Costo de faltante: $8$ por unidad faltante.
\end{itemize}

Se compararon tres politicas:

\begin{center}
\begin{tabular}{lrrl}
\toprule
Politica & $Q$ & $r$ & Interpretacion \\
\midrule
\texttt{bajo\_inventario} & 80 & 45 & Menor inventario promedio, mayor riesgo de faltantes \\
\texttt{balanceada} & 120 & 70 & Punto medio entre mantenimiento y faltantes \\
\texttt{alto\_servicio} & 160 & 95 & Mayor nivel de servicio, mayor mantenimiento \\
\bottomrule
\end{tabular}
\end{center}

En cada dia simulado se aplico la siguiente secuencia:

\begin{enumerate}
    \item Recibir ordenes pendientes cuyo dia de llegada sea el dia actual.
    \item Generar la demanda diaria.
    \item Satisfacer la demanda hasta agotar el inventario disponible.
    \item Registrar unidades faltantes cuando la demanda supera al inventario.
    \item Acumular costo de mantenimiento y costo de faltante.
    \item Calcular la posicion de inventario.
    \item Emitir una orden de cantidad $Q$ si la posicion es menor o igual que $r$.
    \item Acumular el costo fijo de orden y registrar el estado diario.
\end{enumerate}

El tratamiento de faltantes fue de ventas perdidas. Por lo tanto, una unidad no satisfecha se contabiliza como faltante y genera costo, pero no se arrastra como pedido pendiente al dia siguiente. Esta decision mantiene el modelo simple y consistente con la comparacion de politicas de servicio.

Las salidas de Python se guardaron en \texttt{results/inventory/experiments/}. El archivo \texttt{inventario\_corridas.csv} contiene una fila por replica, \texttt{inventario\_resumen.csv} resume cada politica y \texttt{inventario\_series.csv} conserva la evolucion diaria usada para los graficos.

